\section{Discussion}
\label{sec:discussion}


%\todo[inline]{stems from work on \textsc{BiFluX}; comparison with putlenses}
During the design and implementation of \textsc{BiFluX},
it turned out that the putback-based design mixed with regular expression types for XML 
makes the well-behavedness hard to prove
and it also hinders the core idea of putback-based BX.
Based on these observations, \BiGUL\ is proposed and implemented originated from the core of \textsc{BiFluX}
which is datatype-generic and fully verified to guarantee the well-behavedness.
Even though Putlenses~\cite{PaHF14} is the first putback-based combinator language, it is lack of formal verification.
\BiGUL\ is the first fully certified putback-based BX language.
%


%\todo[inline]{extensions: for example, iteration, view dependency}
\BiGUL\ is powerful to express a lot of useful BX programs,
and there are extensions can be made to make it support many other useful features.
%
% Currently alignment only handles the case that
% elements with the same matching key in both source and view are unique % TODO: matching key ?
% which indicates the update operations after alignment focus on single source/view element.
% BiFlux~\cite{PaZH14} also only supports this case,
% while there are situations that more than one elements have the same matching key value,
% the result would be a list of elements for the same key value.
% \BiGUL\ can be extended with iteration operation on list to hanle this case.
% %
% % source dependency rearrangement already done.
% % view dependency need new feature to hanle it.
% Another one is view dependency that means two view elements are denpendent.
% For example the catalog of a book dependends on the strcuture and contents of the book.
% If the title of contents changes, the catalog shall be changed.
% This can be captured by expressing the depdenency between the catalog and the titles in the content.
% Only one of them (e.g. titles in the content) needs to be used to update the corresponding source element.
% From the $get$ perspective, catalog will be recovered through the view dependency.

%\todo[inline]{Haskell implementation}
\BiGUL\ has been fully implemented using Haskell that totally follows the formalization in \Agda\
and tested with several nontrivial examples which are from the BiFluX~\cite{PaZH14} paper.
\todo{implement more examples ?}

\todo[inline]{recursion (the lack thereof)}

\todo[inline]{Towards a dependently typed version}


\todo[inline]{Theory construction in terms of \Agda\ programming\\
Regarding formalisation: ``[W]e are interested in extending what can be calculated precisely because we are not interested in the calculations themselves. Developing techniques that allow more of a derivation to be calculated allows attention to be focussed on the creative parts, which cannot be calculated.''~\cite{Gibbons-BMF}\\
\cite{UFP-HoTT}}

